\section{Limitations and possible improvements of the TCV shape controller}
\label{sec:tcvshapectrl_limitations}
The current design of the TCV shape controller has been quite the challenge for the development team and its integration has turned out to be a delicate matter. Nonetheless, the result is a system that works: given a reference plasma shape and after some tuning, the TCV shape controller is able to control the plasma shape within the limits of the actuators and the diagnostics. Recalling the discussion of the redundant coil configuration of TCV in \Cref{sec:tcv}, which constitutes probably its most peculiar feature, the question arises whether the current shape controller exploits this redundancy to improve the performance of the coil current control or not. Up to this point, it does not.

The PID formulation of the shape controller, described in \cite{mele_design_2024}, does not account for the actuator redundancy offered by the TCV model, which may in turn be exploited to improve the performance of the shape controller itself. Performance indicators to be considered in this context are the overall power consumption of the coils, as well as the saturation limits of the power supplies. In practical experiments, these have turned out to be important aspects to consider when working with TCV.

The above mentioned question leads one to think whether something can be done to improve the present shape control scheme. This is the starting point for the novel work presented in this chapter. An immediate solution is to include the performance optimization goals in the design of the shape controller. This approach, although ideally feasible, leads to the definition of a control problem that is much different from the original one. Determining a solution to it, considering also all the intricacies and complications that are involved in the control of a plasma shape on TCV, is a challenging task. Another approach consists in keeping the current shape controller design, which works and has been validated through numerous experiments, essentially as it is. However, a new control problem can be formulated at this point, considering the current shape controller as part of the dynamic system and with the goal of optimizing the performance of the shape controller by altering its output before it is fed to the plant. This is the approach that has been followed in this work, relying on a specific subject of control theory that is receiving increasing interest in recent years: that of \emph{dynamic control allocation}.

Dynamic control allocation techniques are used to distribute the control effort among a set of redundant actuators in order to optimize a given performance index, following an \emph{allocation policy}. What distinguishes \emph{dynamic} control allocation from similar techniques is that, in the former, the allocation policy is a control law itself that is determined through a control design process and evaluated in real-time. Subsystems that are typically involved in a dynamic control allocation scheme are denoted as \emph{dynamic allocators} or simply as \emph{allocators}. The peculiarity of dynamic allocators is that they are designed as \emph{add-ons} to a preexisting control system, \emph{i.e.}, their design does not require a complete redesign of the control system, which may yield significant complications when optimization objectives are included in the problem. Instead, given the models of the plant and the control system, their design can be carried out \emph{a posteriori} and they can subsequently be integrated into the control scheme as additional subsystems. The art of dynamic control allocation, therefore, lies in the design of these additional components with the goal in mind of not only optimizing the control action in terms of some criteria of practical relevance but also of ensuring that the original control objectives are still met without any alteration of the desired output behavior.

Regarding TCV, a methodological study has been carried out in \cite{tenaglia_dynamic_2024}, in which an early implementation of a dynamic coil current allocator has been integrated in the MATLAB simulation code suite. Its performance has been validated by means of simulations on real shot data, which showed the feasibility of this approach. The rest of this chapter presents the work that has been carried out in this direction. First, a review of the current state of the art about dynamic control allocation is presented, together with a set of formulations that are of interest for the TCV case. Each of these has been considered, leading to some theoretical results that are illustrated in the following. Then, the design of a novel dynamic control allocator for TCV is presented, which is based on the latest advancements in the subject. Experimental results are discussed at the end of the chapter, showing the performance of the proposed novel allocator in a real scenario and constituting the first working implementation of its kind in the literature.
